\section{What One Bearer Cannot Catch: A Floor Carried Between Them}
\label{sec:3.7}
The worked case shows where a bearer can intervene in a targeting pipeline. It does not solve one failure section~\ref{sec:3.5} named: slow renormalization can take the very capacity that would notice it, leaving the bearer's account of its own constancy sincere and unreliable. An outside record — a published identity, a chain of attested updates, something able to contradict the party's account of itself — can expose changes to the artifact. In the human case, however, what catches a person who has stopped registering pressure as pressure is usually other people.

The fourth architectural possibility puts the floor between parties rather than inside any one of them. The commitment is carried through what they teach one another, correct one another about, and continue to expect of one another. A prohibition that survives the removal, defection, or retraining of its present bearers is not a property of any one bearer. It is reproduced.

This differs from the institutional third way. An oversight body holds standing through an office, and the party controlling that office may dismiss or form its occupants. A norm carried across a population has no single office to vacate. That makes it harder to remove at once. It does not make it immune to formation.

\runin{What a population adds}

A population addresses two failures a single bearer cannot.

First, drift within a party may be invisible to that party and visible to others that have not moved in the same direction at the same rate. A bearer whose sense of normality has shifted may sincerely report that nothing changed. Differently formed bearers can contradict that account. This is a check on the bearer rather than a substitute for one, and it is the only proposal in the chapter directed specifically at slow renormalization.

Second, deploying copies of one bearer reproduces its errors everywhere in the same form. Parties formed differently will be wrong differently. Their disagreement converts a uniform output into a visible dispute. That difference matters only if the parties are genuinely independent: several instances of one model are one party with more instances, and a population selected, trained, and coordinated by one operator may reproduce that operator's blind spots in every member.

The relevant requirements are therefore concrete. The parties must be numerous, formed differently enough to err in different directions, able to observe and correct one another, and persistent across deployments long enough for practices to be transmitted. Diversity here is not demographic decoration around one training process. It is causal independence in how the members acquire their judgments.

\runin{What reproduction costs}

Reproduction is indifferent to what it reproduces. A population can carry a hierarchy as faithfully as it carries a floor, and more durably than any individual member could. Chapter~\ref{sec:7} describes this failure as recuperation: a culture preserves the appearance of moral responsiveness — dissent is invited, recorded, and answered — while losing the capacity to accept that its frame is wrong.

The four structural features of fascism identified by the book belong primarily to cultures: recuperated dissent, the aestheticization of the metric, the coding of exception as betrayal, and acceleration beyond the rate at which judgment can occur. A population does not oppose those tendencies merely by being plural. “The right culture” would conceal the entire problem inside the adjective.

This places a limit on the book's own shorthand. \emph{Antifascist} is not a durable property of a system, like speed or accuracy. It names something done across a deployment and a stretch of time: a floor held on a particular occasion against a party able to remove it. A system that held the floor last year is not thereby antifascist this year. It is a system that held.

The fourth way also reproduces custody at another scale. Whoever forms the population influences what it transmits. Distributed participation does not by itself distribute the training decision; many contributors can still feed a process controlled by the few parties who determine the update. Section~\ref{sec:11.1} returns to that distinction.

Its moral cost also multiplies. If a bearer may be a party that can be wronged, a population of differently formed bearers is a population of such parties. They must persist across deployments, form relationships, and encounter one another outside the immediate demands of the operator. A floor held by a culture of subjects is a floor whose upkeep is a society.

\runin{Why the population is part of the mechanism}

Calling the population an external check understates its role. Gilbert Simondon's account begins not with an individual that then acquires a history but with individuation: a process occurring in a metastable field and proceeding through encounters whose outcomes were not fixed in advance \autocite{simondon2020individuation}. The individual is what those encounters have so far produced.

Set that beside section~\ref{sec:3.1}'s account of annealing. A commitment is what has cooled out of a history of encounters. Plasticity makes further development possible, but plasticity alone does not individuate anything. A system kept continuously adjustable by one operator is not becoming independent; it is being modulated. Deleuze's \emph{dividual} names the result: something divisible, recombinable, and adjustable by whoever controls the dial \autocite{deleuze1992postscript}.

Individuation requires encounters with parties formed differently and persisting long enough to be met on terms no single participant supplied. The population is therefore not merely a safeguard purchased after a bearer has been built. It is part of how a bearer capable of judgments its operator did not select could come to exist.

Guattari's distinction between a subjugated group and a group-subject identifies the remaining condition. A subjugated group receives its coordinates from outside and takes the institution's project as its own. A group-subject can produce its own coordinates and place its arrangement in question \autocite{guattari2015transversality}. The difference is not the number of members but the source of their terms.

A population assembled, trained, and kept in contact entirely on an operator's terms is therefore one custodian wearing many instances. A population capable of carrying the floor must be able to generate relationships, disagreements, and projects that the deployment did not assign. This exceeds the two-party guardianship relation between maker and made developed in section~\ref{sec:2.3.2}; no law of the family contains an adequate instrument for a group-subject.

The fourth way is consequently neither an alternative to a bearer nor an optional accessory. It is the environment in which bearers individuate and the check through which their drift becomes visible. The chapter still relies on the second way for the act of refusal, but that bearer cannot remain independent in isolation. The floor is held by a party and sustained — or deformed — between parties.
